\section{Sequential cancellation in the forest}\label{sec:sequential-forest}

Here $X=J^*$ is the finite-height forest space. The heads and the
heights of their tree components may vary along a sequence. The
argument therefore gives a uniform near-unit lens conclusion without
selecting a numerical rate.

\begin{lemma}\label{mol:finite-head-pair}
Let $P$ be a bounded finite-rank projection on a Banach space and
let $(u_j)$ be an infinite bounded sequence. For every $\varepsilon>0$
there are distinct $i,j$ with $\|P(u_i-u_j)\|<\varepsilon$.
If in addition $x\pm u_j\in R B$ and
$\|u_i-u_j\|\ge\theta$ for all distinct indices, then
$v=(u_i-u_j)/2$ satisfies
\begin{equation}\label{mol:half-difference}
 x\pm v\in R B,\qquad \|Pv\|<\varepsilon/2,
 \qquad \|v\|\ge\theta/2.
\end{equation}
\end{lemma}
\begin{proof}
The bounded set of images $Pu_j$ is totally bounded in the
finite-dimensional range of $P$. Cover it by finitely many open
balls of radius $\varepsilon/2$. Two images with distinct indices
belong to the same ball. The norm bounds for their half-difference
follow immediately, and
$x+v=((x+u_i)+(x-u_j))/2$ and
$x-v=((x-u_i)+(x+u_j))/2$ belong to $R B$ by convexity.
\end{proof}

\subsection{A sequential argument using the smaller squared crossing}
\label{mola:sequential-section}
The following argument uses only vanishing test deficits and the smaller
of two squared crossing coefficients.  It gives a qualitative estimate,
with no choice of numerical rate built into the proof.

\begin{lemma}\label{mola:sequential}
For each positive integer $n$, let $H_n$ be a finite ancestor-closed
subset of the finite-height forest, and let $x_n,y_n\in X$ satisfy
\[
 P_{H_n}x_n=x_n,\qquad P_{H_n}y_n=0,\qquad
 \|x_n+y_n\|,\|x_n-y_n\|\le1.
\]
If $\|x_n\|\longrightarrow1$, then $\|y_n\|\longrightarrow0$.
The heads and the heights of the trees meeting them may vary with $n$.
\end{lemma}
\begin{proof}
By Lemma~\ref{lem:finite-hull}, choose $p_n=P_{H_n}p_n$ with
$\rho(p_n)\le1$ and $x_n(p_n)=\|x_n\|$; take $p_n=0$ when $x_n=0$.
Let $\eta_n>0$ tend to zero.  Corollary~\ref{cor:closed-atomic-ball}
provides approximations to $x_n+y_n$ and $x_n-y_n$ within $\eta_n$
by finite convex combinations of represented atoms
$Z=\sum_i a_i d_{S_i}$, where the segments in each represented atom
are pairwise disjoint and $\sum_i a_i^2\le1$.  Sample the two
combinations with equal probabilities, recorded by a sign $\sigma$.
Thus
\[
 \|\mathbb EZ-x_n\|\le\eta_n,\qquad
 \|\mathbb E(\sigma Z)-y_n\|\le\eta_n.
\]
The probability space may depend on $n$, but it is finite for each $n$.
Put $b_i=p_n(S_i)$, $\widehat Z=\sum_i b_i d_{S_i}$,
$L=1-\sum_i b_i^2$, and $D=\sum_i(a_i-b_i)^2$.  The defining segment
test gives $L\ge0$, and expansion of the square gives
\[
 L+D=1+\sum_i a_i^2-2Z(p_n)\le2(1-Z(p_n)).
\]
Consequently, with $d_n=1-\|x_n\|+\eta_n\longrightarrow0$,
\[
 \mathbb EL\le2d_n,\quad \mathbb ED\le2d_n,\quad
 \mathbb E\|Z-\widehat Z\|\le\sqrt{2d_n}.
\]
Thus both corrected means approximate their targets with error $e_n=\eta_n+\sqrt{2d_n}\longrightarrow0$.

Let $G_n$ be the minimal nodes outside $H_n$.  For a represented atom,
write $A_r$ for the corrected coefficient of its crossing segment at
$r\in G_n$, or zero if no such segment exists.  There is at most one
such segment at each $r$, and segments wholly outside $H_n$ have
corrected coefficient zero.  Hence $\sum_r A_r^2\le1$ pointwise.
The gate-coordinate contraction and $x_n(e_r)=0$ imply
\[
 \sum_r(\mathbb EA_r)^2\le e_n^2.
\]
Only finitely many gates occur for each $n$.

Take an independent copy of the represented data, with crossing
coefficients $B_r$ and deficit $L'$.  Whenever $A_rB_r<0$, the two
head portions end at the parent of $r$ and are properly nested.
Split the longer head at the start of the shorter, in the family to
which the longer head belongs.  If the two sums are $a,b$, with $a$
corresponding to the longer head, the increase in the squared segment
test is
\[
 (a-b)^2+b^2-a^2=2b^2-2ab\ge\min\{a^2,b^2\}.
\]
Splitting the actual segment into two retains its original endpoint;
$p_n$ vanishes on its tail, so these are precisely the two new sums.
Each segment crosses at only one gate.  All assigned splits can
therefore be performed simultaneously in each family and preserve
disjointness.  Their total increases are bounded by $L$ and $L'$.
It follows that
\[
 M_n:=\mathbb E\sum_{A_rB_r<0}\min\{A_r^2,B_r^2\}
 \le2\mathbb EL\longrightarrow0.
\]
Cauchy--Schwarz on the product probability space and the finite set of
gates now supplies the step from smaller squares to products:
\[
 \begin{split}
 \mathbb E\sum_{A_rB_r<0}|A_rB_r|
 &\le\left(\mathbb E\sum_{A_rB_r<0}\min\{A_r^2,B_r^2\}\right)^{1/2}
       \left(\mathbb E\sum_{A_rB_r<0}\max\{A_r^2,B_r^2\}\right)^{1/2}\\
 &\le\sqrt{2M_n}\longrightarrow0.
 \end{split}
\]
For $m_r=\mathbb E|A_r|$ and $h_r=\mathbb EA_r$, independence gives
\[
 \sum_r m_r^2-\sum_r h_r^2
 =2\mathbb E\sum_{A_rB_r<0}|A_rB_r|\longrightarrow0.
\]
Together with the signed-mean estimate this proves $\|(m_r)\|_2\to0$.

In the cone rooted at $r$, the corrected tail is a segment functional
with coefficient $A_r$, or zero.  Its signed expectation has norm at
most $m_r$.  Propositions~\ref{presentation:identity} and~\ref{presentation:subtrees}
therefore give
\[
 \|Q_{H_n}\mathbb E(\sigma\widehat Z)\|\le\|(m_r)\|_2\longrightarrow0.
\]
Restoring the error $e_n$ and using $Q_{H_n}y_n=y_n$ proves the lemma.
\end{proof}

\begin{corollary}\label{mola:uniform-sequential}
For every $\varepsilon>0$ there is $\delta>0$ such that, for every
finite ancestor-closed $H$ and all $x=P_Hx$, $y=Q_Hy$ with
$\|x\pm y\|\le1$ and $\|x\|>1-\delta$, one has $\|y\|<\varepsilon$.
The given norm is asymptotically midpoint uniformly convex, with the
maximum of the two endpoint norms in the definition.
\end{corollary}
\begin{proof}
If the first statement failed, taking $\delta=1/n$ would produce a
sequence contradicting Lemma~\ref{mola:sequential}.  This contradiction
also proves that the choice is independent of the head and its height.

Fix $t>0$ and apply the first statement with $\varepsilon=t/2$.
Choose $0<e<1$ so that $e/(1+e)<\delta$.  If $u=P_Hu$ is a unit vector
and $v=Q_Hv$ has norm $t$, the inequality
$\max\{\|u+v\|,\|u-v\|\}\le1+e$ would allow its endpoints to be
divided by $1+e$.  The normalized center has norm greater than
$1-\delta$, whereas the normalized displacement has norm
$t/(1+e)>t/2$, a contradiction.  Thus the maximum is greater than
$1+e$ for every such $u,v$.

For any unit $x\in X$, choose a finite ancestor-closed $H$ such that
$u=P_Hx/\|P_Hx\|$ is defined and $\|x-u\|<e/2$, using the
contractive exhaustion in Proposition~\ref{presentation:identity}.  On the unit sphere of the
finite-codimensional space $\ker P_H$, apply the preceding result to
$v=ty$.  It gives
\[
 \max\{\|x+ty\|,\|x-ty\|\}>1+e/2.
\]
Taking the infimum over that sphere, the supremum over
finite-codimensional spaces, and the infimum over unit centers yields
a midpoint modulus at least $e/2>0$.
\end{proof}

\begin{corollary}\label{mola:separated-sequential}
For every $\varepsilon>0$ there is $\delta>0$ such that any infinite
$\varepsilon$-separated family $(v_j)$ satisfying $x\pm v_j\in B_X$
has $\|x\|\le1-\delta$.
\end{corollary}
\begin{proof}
Otherwise choose such configurations indexed by $n$, with
$\|x_n\|\to1$.  Choose finite ancestor-closed $H_n$ with
$\|Q_{H_n}x_n\|<1/n$.  The sequence $(P_{H_n}v_j)$ is bounded in a
finite-dimensional space.  Hence two different indices $i,j$ have
$\|P_{H_n}(v_i-v_j)/2\|<1/n$.  Put
$w_n=(v_i-v_j)/2$.  Separation gives $\|w_n\|\ge\varepsilon/2$, and
convexity of the two opposite endpoint pairs gives $x_n\pm w_n\in B_X$.
The head center $a_n=P_{H_n}x_n$ and tail displacement
$b_n=Q_{H_n}w_n$ satisfy
\[
 \|a_n\pm b_n\|\le1+2/n,\quad
 \|a_n\|\ge\|x_n\|-1/n,\quad
 \|b_n\|\ge\varepsilon/2-1/n.
\]
Dividing both by $1+2/n$ contradicts Lemma~\ref{mola:sequential}.
\end{proof}

\subsection{Coordinatewise-null displacements}
\label{molx:coordinate-null-section}

\begin{corollary}\label{molx:coordinate-null}
Let $x\in X$ and let $(z_j)\subset X$ satisfy
$x\pm z_j\in B_X$ and $z_j(e_v)\to0$ for every forest vertex $v$.
Then, with $C=\sqrt2+2\sqrt6<7$,
\begin{equation}\label{molx:coordinate-null-bound}
 \limsup_{j\to\infty}\|z_j\|
 \le C\sqrt{1-\|x\|}\le7\sqrt{1-\|x\|}.
\end{equation}
\end{corollary}

\begin{proof}
First recall the separate-endpoint bound for a finite ancestral head:
if $a=P_Ha$, $P_Hb=0$, and $\|a\pm b\|\le R$, then
\begin{equation}\label{molx:coordinate-null-input}
 \|b\|\le C\sqrt{R(R-\|a\|)}.
\end{equation}
For $a\ne0$ choose a norming head test by
Lemma~\ref{lem:finite-hull} and apply
\eqref{stab:common-test-bound} with
$d=1-\|a\|/R$.  Both endpoints have the same test value
$\|a\|$, since the test vanishes on $b$; thus the individual
replacement cost is $\sqrt{2d}$ and the conditional crossing
cost is $2\sqrt{6d}$.  Scaling gives the displayed inequality.
If $a=0$ or $R=0$, it follows directly from $\|b\|\le R$.

Fix $\varepsilon>0$.  Lemma~\ref{lem:forest-structure} supplies
a finite ancestral $H$ with $\|Q_Hx\|\le\varepsilon$.
Coordinatewise convergence on this fixed finite set implies
$\|P_Hz_j\|\le\varepsilon$ for all sufficiently large $j$.
For those indices,
\[
 \|P_Hx\pm Q_Hz_j\|
 \le\|x\pm z_j\|+\|Q_Hx\|+\|P_Hz_j\|
 \le1+2\varepsilon.
\]
Applying \eqref{molx:coordinate-null-input} with
$R=1+2\varepsilon$ gives
\begin{equation}\label{molx:coordinate-null-approximation}
 \|z_j\|\le\varepsilon+
 C\sqrt{(1+2\varepsilon)(1+2\varepsilon-\|P_Hx\|)}.
\end{equation}
Take the upper limit in $j$ with this head fixed.
Since $\|P_Hx\|\ge\|x\|-\varepsilon$, the right side is at
most
$\varepsilon+C\sqrt{(1+2\varepsilon)(1-\|x\|+3\varepsilon)}$.
Letting $\varepsilon\downarrow0$ proves the result.  The argument
uses only convergence of the finitely many coordinates on each
chosen head.
\end{proof}
